\section{방법}

\subsection{스크린샷 선별과 좌표 정규화}

WebUI-350K \citep{biglab2024webui350k}에서 Tranco 도메인 순위,
도메인별 상한, 라운드로빈 선별을 이용해 영어 데스크톱 스크린샷
150K장을 선택한다. 각
예제는 $1920\times1080$ viewport slice이다. Chrome 접근성 트리에서
요소의 역할, 이름, 클릭 가능 여부, bounding box를 얻는다. 이미지
너비와 높이가 각각 $W,H$이고 bounding box가 $(x,y,w,h)$일 때
중심점을 다음과 같이 공통 정수 좌표계로 변환한다.
\begin{equation}
\hat{x}=\operatorname{clip}\!\left(
\operatorname{round}\left(\frac{x+w/2}{W}\times999\right),0,999\right),
\end{equation}
$\hat{y}$도 같은 방식으로 계산한다. 합성 모델은 계산된 좌표를
프롬프트로 전달받아 그대로 인용하며, 픽셀만 보고 좌표를 추정하지
않는다.

\subsection{2단계 evidence 추출과 verbalization}

Qwen3.5-397B-A17B teacher를 이용해 설명을 두 단계로 합성한다.
\textbf{Stage 1: 구조화 사실 추출.} Stage 1은 스크린샷과 화면에
보이는 접근성 트리 요소를 입력받아
페이지 유형, 팝업 상태, 가시 텍스트 사실, 현재 행동 가능한 요소를
구조화한다. 각 affordance는 이름, 행동 유형, 예상 결과, 좌표를
포함한다. 접근성 트리는 웹 UI의 시각ㆍ텍스트 요소를 합성
supervision으로 변환하는 구조화 근거로 활용되어 왔다
\citep{liu2025multiui}. 총 148K개의 유효한 사실 레코드를 얻었으며
세 설명 조건 모두 이 파일을 공유한다. 따라서 조건 간 차이는 teacher가
서로 다른 사실을 추출했기 때문이 아니라, 같은 사실을 어떤
정보 budget으로 언어화했기 때문에 발생한다.

\textbf{Stage 2: 포괄성 수준별 설명 생성.} Stage 2는 같은
스크린샷과 Stage 1 출력을 입력받아 Concise, Balanced 또는
Comprehensive 설명을 생성한다. 이처럼
시각 정보의 추출 단계를 공유한 뒤 언어화 단계만 분기함으로써,
비교 대상이 사실 추출 품질이 아니라 target의 information
budget과 evidence coverage가 되도록 한다.

\paragraph{Concise.}
Concise prompt는 3--5문장, 약 90--130단어로 페이지 유형, 주된 목적,
가장 중요한 가시 상태, 최대 두 개의 핵심 상호작용 목표를
서술하도록 한다. 출력에 행동 요소가 등장하면 제공된 좌표를 반드시
함께 보존한다. 즉 화면 전체를 열거하지 않고 salient evidence를
선택하고 압축한다. 최종 데이터는 143,097건이며 평균
110.2단어, 4.24문장이다.

\paragraph{Balanced.}
Balanced prompt는 6--9문장, 약 170--230단어를 목표로 한다.
페이지 목적과 현재 상태를 유지하면서 각 주요 section에서 대표
사실을 선택하고, primary task flow를 반드시 포함한 뒤 서로 다른
기능을 대표하는 interaction을 추가한다. interaction goal은
3--4개, actionable element는 최대 6개로 제한하며, 언급한 모든
행동 요소에는 Stage 1 좌표를 붙인다. 표와 반복 목록은 header와
대표 값만 포함하고, 선택한 핵심 interaction 1--2개의 expected
결과만 추측형으로 기술한다. Concise의 모든 항목을 sample마다
복사하도록 강제하지는 않지만, primary task를 부차적인 navigation
요소보다 우선하여 평균적으로 중간 수준의 coverage를 형성한다.

\paragraph{Comprehensive.}
Comprehensive prompt는 제공된 가시 사실과 affordance를 하나의
자기완결적 장문에 거의 모두 포함하도록 요구한다. 표의 행과 열,
요소 좌표를 보존하고, 관찰된 사실과 아직 실행하지 않은 행동의
예상 결과를 구분한다. 이미지 해시에 따라 전개 순서와 문체를
결정론적으로 변주하므로, 이 조건은 높은 coverage와 함께 다양한
information ordering을 갖는 실용적 recipe다. 반복 퇴화가 발생한
25건을 제거한 뒤 147,967개의 설명을 얻었으며 평균 길이는
3,656자이다.

세 조건은 동일 screenshot 수를 목표로 하지만 target-token 수는
의도적으로 다르다. 따라서 주 비교는 token-matched length
ablation이 아니라 실제 verbalization recipe의 information
completeness 비교이다. 최종 비교에는 세 설명이 모두 유효한 동일
screenshot의 교집합을 사용하고 sample 수와 update 수를 맞춘다.
추가로 Grounding-only 조건을 두어 description이 없는 공통
grounding mixture와 비교한다.

\begin{table}[t]
\centering
\small
\setlength{\tabcolsep}{5pt}
\begin{tabular}{lrr}
\toprule
 & Short & Long \\
\midrule
화면 설명 & 143K & 148K \\
공통 그라운딩 & 1.2M & 1.2M \\
전체 PT 예제 & 1.37M & 1.38M \\
\bottomrule
\end{tabular}
\caption{\textbf{사전학습 데이터 구성.} 두 조건은 같은 그라운딩
데이터와 좌표계를 공유한다. 표는 전체 구축 규모이며, 통제 비교는
두 설명이 모두 유효한 동일 화면의 교집합에서 수행한다.}
\label{tab:data}
\end{table}


\paragraph{프로그램 기반 검증.}
모든 조건에서 생성 좌표가 Stage 1 좌표 집합의 부분집합인지
검사한다. Concise와 Balanced에는 길이ㆍ문장 수ㆍactionable-element
budget 검사를 추가하고, Comprehensive에는 반복 퇴화 필터를
적용한다. 현재 완료된 Concise 생성 결과의 좌표를 대조한 결과,
약 148K개 입력 가운데 두 차례 재시도 후에도 약 4.9K건이 규칙을
위반하였다. 이 중 약 4.7K건은 존재하지 않는 좌표를 생성했고,
나머지는 문장 수 또는 단어 수 제약을 어겼다. 최종 집합에는 약
392K개의 좌표 언급(예제당 평균 2.74개)이 있으며 검출된 좌표
발명은 0건이다. 이 결과는 합성 데이터의 구조적 제약을 프롬프트에만
맡기지 않고 검증기로 강제해야 함을 보여준다.
이는 합성 UI 기능 주석에 rejection과 verification을 적용한
AutoGUI \citep{li2025autogui}, 그리고 존재하지 않는 요소에 좌표를
생성하는 fabricated hallucination을 보고한 GUI-HalluBench
\citep{zhang2026guihallu}와 일관된다.

\subsection{공통 그라운딩과 비교 조건}

각 설명 집합을 동일한 1.2M개의 영어 UGround 예제와 혼합한다
\citep{gou2025uground}. UGround는 GUI 요소를 가리키는 표현과
해당 요소의 위치를 연결하는 시각 그라운딩 supervision을 제공한다.
본 연구에서는 이를 동일한 0--999 좌표계의 한 점으로 표현한다.
세 혼합 데이터는 seed 42와 동일한
\texttt{prompt}/\texttt{completion}/\texttt{images} 스키마를 사용한다.
Grounding-only는 이 1.2M 예제만 학습한다. 따라서
Grounding-only--Concise는 선택적 description의 추가 효과,
Grounding-only--Comprehensive는 포괄적 description의 추가 효과,
Concise--Balanced--Comprehensive는 증가하는 information budget의
효과를 각각 보여준다. 더 포괄적인 조건이 description 안에서 많은
좌표를 언급하더라도 별도의 UGround 예제를 더 받지는 않는다.
